%% ECON 282E -- Session 3, Deck B1: Your Toolkit
%% Budget: 16 frames (~22 minutes). Provenance: slides/SOURCES.md.
%% Ported from QM Lec. 3 sections 1-4 (tech stack, installation, data types),
%% with three frames written new: environments, Git, and Colab, which the source
%% does not cover at all.

%% ECON 282E -- Foundations of Macroeconomics (UC Riverside, Fall 2026)
%% Shared Beamer preamble for every deck in slides/S01 ... slides/S10.
%%
%% Usage (a deck lives in slides/S0X/ and is compiled from that folder):
%%   %% ECON 282E -- Foundations of Macroeconomics (UC Riverside, Fall 2026)
%% Shared Beamer preamble for every deck in slides/S01 ... slides/S10.
%%
%% Usage (a deck lives in slides/S0X/ and is compiled from that folder):
%%   %% ECON 282E -- Foundations of Macroeconomics (UC Riverside, Fall 2026)
%% Shared Beamer preamble for every deck in slides/S01 ... slides/S10.
%%
%% Usage (a deck lives in slides/S0X/ and is compiled from that folder):
%%   \input{../preamble/course_preamble.tex}
%%   \decktitle{1}{A1}{Who I Am \& Why This Course}{September 24, 2026}
%%   \begin{document}
%%   \makedecktitle
%%   ...
%%
%% Adapted from Courses/AI-ECON-2026/source/shared_preamble.tex (Metropolis theme,
%% concept/caution/example/takeaway boxes, \sectiondivider). Changes for this course:
%%   * course title block (\decktitle / \makedecktitle) instead of \makebeamertitle
%%   * \zh{...} typesets Chinese (book titles) under pdflatex via CJKutf8
%%   * researchbox for the "Research in the wild" frames
%%   * section pages off; TikZ libraries and the course map loaded here
%% Build with pdflatex only (two passes); tools/build_slides.py does this for you.

\documentclass[10pt,english,aspectratio=169]{beamer}

\usepackage[T1]{fontenc}
\usepackage{textcomp}
\usepackage[utf8]{inputenc}
\usepackage{babel}
\usepackage{CJKutf8}

\usepackage{amsmath, amssymb, amsthm, graphicx}
\usepackage{booktabs, multirow, tabularx}
\usepackage{tikz}
\usetikzlibrary{positioning,arrows.meta,shapes,calc,fit,backgrounds}
\usepackage{pgfplots}
\pgfplotsset{compat=1.16}
\usepackage[most]{tcolorbox}
\tcbuselibrary{skins, breakable}
\usepackage{listings}
\usepackage{xcolor}

\lstset{
  basicstyle=\ttfamily\small,
  keywordstyle=\color{blue!80!black}\bfseries,
  commentstyle=\color{green!50!black}\itshape,
  stringstyle=\color{orange!80!black},
  backgroundcolor=\color{gray!8},
  frame=single,
  rulecolor=\color{gray!40},
  breaklines=true,
  showstringspaces=false,
  numbers=none,
}

\usetheme{metropolis}
\usefonttheme{professionalfonts}
\metroset{sectionpage=none}

\definecolor{LightGray}{RGB}{230,230,230}
\definecolor{RedTitle}{RGB}{0,0,200}
\definecolor{VeryLightGray}{RGB}{245,245,245}
\definecolor{DarkText}{RGB}{0,0,0}
\definecolor{UCRBlue}{RGB}{0,45,106}
\definecolor{UCRGold}{RGB}{241,171,0}

% Stage colours of the course arc (used by the course map and elsewhere)
\colorlet{StageTrunk}{blue!12}
\colorlet{StageBranch}{cyan!14}
\colorlet{StageMachine}{gray!18}
\colorlet{StageWall}{red!14}
\colorlet{StageTools}{orange!20}
\colorlet{StageData}{green!16}
\colorlet{StageAgents}{violet!14}

\hypersetup{bookmarks=false,breaklinks=true,pdfborder={0 0 0},colorlinks=true,
  linkcolor=DarkText,urlcolor=RedTitle}

\setbeamercolor{background canvas}{bg=white}
\setbeamercolor{title}{fg=RedTitle}
\setbeamercolor{frametitle}{bg=VeryLightGray,fg=DarkText}
\setbeamercolor{normal text}{fg=DarkText}
\setbeamercolor{alerted text}{fg=RedTitle}
\setbeamersize{text margin left=5mm, text margin right=5mm}
\addtobeamertemplate{frametitle}{}{\vspace{-0.5em}}

\setbeamertemplate{navigation symbols}{}
\setbeamertemplate{footline}{%
  \hspace{5mm}{\usebeamerfont{page number in head/foot}\color{black!45}%
  ECON 282E \,$\cdot$\, \insertshortdeck}\hfill%
  \usebeamercolor[fg]{page number in head/foot}%
  \usebeamerfont{page number in head/foot}%
  \insertframenumber\,/\,\inserttotalframenumber\kern1em\vskip2pt%
}

% ---------------------------------------------------------------------------
% Chinese text under pdflatex (book titles etc.)
\newcommand{\zh}[1]{\begin{CJK*}{UTF8}{gbsn}#1\end{CJK*}}

% ---------------------------------------------------------------------------
% Deck title block
%   \decktitle{session}{deck id}{title}{date}
\newcommand{\insertshortdeck}{}
\newcommand{\decksession}{}
\newcommand{\deckid}{}
\newcommand{\decktitletext}{}
\newcommand{\deckdate}{}
\newcommand{\decktitle}[4]{%
  \renewcommand{\decksession}{#1}%
  \renewcommand{\deckid}{#2}%
  \renewcommand{\decktitletext}{#3}%
  \renewcommand{\deckdate}{#4}%
  \renewcommand{\insertshortdeck}{Session #1 \,$\cdot$\, Deck #1.#2}%
  \title{#3}%
  \author{Zhigang Feng}%
  \date{#4}%
}
\newcommand{\makedecktitle}{%
  \begin{frame}[plain,noframenumbering]
    \begin{tikzpicture}[remember picture,overlay]
      \fill[UCRBlue] (current page.north west) rectangle ([yshift=-0.55cm]current page.north east);
      \fill[UCRGold] ([yshift=-0.55cm]current page.north west) rectangle ([yshift=-0.62cm]current page.north east);
      \node[anchor=west, text=white, font=\small\bfseries] at ([xshift=5mm,yshift=-0.275cm]current page.north west)
        {ECON 282E \,$\cdot$\, Foundations of Macroeconomics};
      \node[anchor=east, text=white, font=\small] at ([xshift=-5mm,yshift=-0.275cm]current page.north east)
        {UC Riverside \,$\cdot$\, Fall 2026};
    \end{tikzpicture}
    \vspace*{1.1cm}
    {\usebeamerfont{title}\color{black!55}\large Session \decksession\ \,$\cdot$\, Deck \decksession.\deckid\par}
    \vspace{0.25cm}
    {\usebeamerfont{title}\color{RedTitle}\LARGE\bfseries \decktitletext\par}
    \vspace{0.7cm}
    {\large Zhigang Feng\par}
    \vspace{0.15cm}
    {\color{black!65}\deckdate\par}
    \vfill
    {\footnotesize\itshape\color{black!55}Build it on paper $\;\to\;$ solve it on the machine $\;\to\;$ push it to the frontier with AI\par}
  \end{frame}}

% ---------------------------------------------------------------------------
% Section divider (plain frame, not counted)
\newcommand{\sectiondivider}[2]{%
\begin{frame}[plain,noframenumbering]
\begin{center}
\vfill
{\Large\textbf{#1}\par}
\vspace{0.35cm}
{\normalsize #2\par}
\vfill
\end{center}
\end{frame}
}

% ---------------------------------------------------------------------------
% Boxes
\newtcolorbox{conceptbox}[1][]{colback=blue!3, colframe=RedTitle!55, boxrule=0.35mm,
  arc=1mm, left=2mm, right=2mm, top=1mm, bottom=1mm, #1}
\newtcolorbox{cautionbox}[1][]{colback=red!3, colframe=red!50!black, boxrule=0.35mm,
  arc=1mm, left=2mm, right=2mm, top=1mm, bottom=1mm, #1}
\newtcolorbox{examplebox}[1][]{colback=green!3, colframe=green!45!black, boxrule=0.35mm,
  arc=1mm, left=2mm, right=2mm, top=1mm, bottom=1mm, #1}
\newtcolorbox{takeawaybox}[1][]{colback=VeryLightGray, colframe=RedTitle!55, boxrule=0.35mm,
  arc=1mm, left=2mm, right=2mm, top=1mm, bottom=1mm, #1}
\newtcolorbox{researchbox}[1][]{enhanced, colback=violet!4, colframe=violet!60!black,
  boxrule=0.35mm, arc=1mm, left=2mm, right=2mm, top=1mm, bottom=1mm,
  fonttitle=\bfseries\small, title={Research in the wild}, #1}

\newcommand{\compacttables}{\renewcommand{\arraystretch}{1.15}}
\newcommand{\src}[1]{{\tiny\color{black!55}#1}}

% ---------------------------------------------------------------------------
\input{../preamble/course_map.tex}

%%   \decktitle{1}{A1}{Who I Am \& Why This Course}{September 24, 2026}
%%   \begin{document}
%%   \makedecktitle
%%   ...
%%
%% Adapted from Courses/AI-ECON-2026/source/shared_preamble.tex (Metropolis theme,
%% concept/caution/example/takeaway boxes, \sectiondivider). Changes for this course:
%%   * course title block (\decktitle / \makedecktitle) instead of \makebeamertitle
%%   * \zh{...} typesets Chinese (book titles) under pdflatex via CJKutf8
%%   * researchbox for the "Research in the wild" frames
%%   * section pages off; TikZ libraries and the course map loaded here
%% Build with pdflatex only (two passes); tools/build_slides.py does this for you.

\documentclass[10pt,english,aspectratio=169]{beamer}

\usepackage[T1]{fontenc}
\usepackage{textcomp}
\usepackage[utf8]{inputenc}
\usepackage{babel}
\usepackage{CJKutf8}

\usepackage{amsmath, amssymb, amsthm, graphicx}
\usepackage{booktabs, multirow, tabularx}
\usepackage{tikz}
\usetikzlibrary{positioning,arrows.meta,shapes,calc,fit,backgrounds}
\usepackage{pgfplots}
\pgfplotsset{compat=1.16}
\usepackage[most]{tcolorbox}
\tcbuselibrary{skins, breakable}
\usepackage{listings}
\usepackage{xcolor}

\lstset{
  basicstyle=\ttfamily\small,
  keywordstyle=\color{blue!80!black}\bfseries,
  commentstyle=\color{green!50!black}\itshape,
  stringstyle=\color{orange!80!black},
  backgroundcolor=\color{gray!8},
  frame=single,
  rulecolor=\color{gray!40},
  breaklines=true,
  showstringspaces=false,
  numbers=none,
}

\usetheme{metropolis}
\usefonttheme{professionalfonts}
\metroset{sectionpage=none}

\definecolor{LightGray}{RGB}{230,230,230}
\definecolor{RedTitle}{RGB}{0,0,200}
\definecolor{VeryLightGray}{RGB}{245,245,245}
\definecolor{DarkText}{RGB}{0,0,0}
\definecolor{UCRBlue}{RGB}{0,45,106}
\definecolor{UCRGold}{RGB}{241,171,0}

% Stage colours of the course arc (used by the course map and elsewhere)
\colorlet{StageTrunk}{blue!12}
\colorlet{StageBranch}{cyan!14}
\colorlet{StageMachine}{gray!18}
\colorlet{StageWall}{red!14}
\colorlet{StageTools}{orange!20}
\colorlet{StageData}{green!16}
\colorlet{StageAgents}{violet!14}

\hypersetup{bookmarks=false,breaklinks=true,pdfborder={0 0 0},colorlinks=true,
  linkcolor=DarkText,urlcolor=RedTitle}

\setbeamercolor{background canvas}{bg=white}
\setbeamercolor{title}{fg=RedTitle}
\setbeamercolor{frametitle}{bg=VeryLightGray,fg=DarkText}
\setbeamercolor{normal text}{fg=DarkText}
\setbeamercolor{alerted text}{fg=RedTitle}
\setbeamersize{text margin left=5mm, text margin right=5mm}
\addtobeamertemplate{frametitle}{}{\vspace{-0.5em}}

\setbeamertemplate{navigation symbols}{}
\setbeamertemplate{footline}{%
  \hspace{5mm}{\usebeamerfont{page number in head/foot}\color{black!45}%
  ECON 282E \,$\cdot$\, \insertshortdeck}\hfill%
  \usebeamercolor[fg]{page number in head/foot}%
  \usebeamerfont{page number in head/foot}%
  \insertframenumber\,/\,\inserttotalframenumber\kern1em\vskip2pt%
}

% ---------------------------------------------------------------------------
% Chinese text under pdflatex (book titles etc.)
\newcommand{\zh}[1]{\begin{CJK*}{UTF8}{gbsn}#1\end{CJK*}}

% ---------------------------------------------------------------------------
% Deck title block
%   \decktitle{session}{deck id}{title}{date}
\newcommand{\insertshortdeck}{}
\newcommand{\decksession}{}
\newcommand{\deckid}{}
\newcommand{\decktitletext}{}
\newcommand{\deckdate}{}
\newcommand{\decktitle}[4]{%
  \renewcommand{\decksession}{#1}%
  \renewcommand{\deckid}{#2}%
  \renewcommand{\decktitletext}{#3}%
  \renewcommand{\deckdate}{#4}%
  \renewcommand{\insertshortdeck}{Session #1 \,$\cdot$\, Deck #1.#2}%
  \title{#3}%
  \author{Zhigang Feng}%
  \date{#4}%
}
\newcommand{\makedecktitle}{%
  \begin{frame}[plain,noframenumbering]
    \begin{tikzpicture}[remember picture,overlay]
      \fill[UCRBlue] (current page.north west) rectangle ([yshift=-0.55cm]current page.north east);
      \fill[UCRGold] ([yshift=-0.55cm]current page.north west) rectangle ([yshift=-0.62cm]current page.north east);
      \node[anchor=west, text=white, font=\small\bfseries] at ([xshift=5mm,yshift=-0.275cm]current page.north west)
        {ECON 282E \,$\cdot$\, Foundations of Macroeconomics};
      \node[anchor=east, text=white, font=\small] at ([xshift=-5mm,yshift=-0.275cm]current page.north east)
        {UC Riverside \,$\cdot$\, Fall 2026};
    \end{tikzpicture}
    \vspace*{1.1cm}
    {\usebeamerfont{title}\color{black!55}\large Session \decksession\ \,$\cdot$\, Deck \decksession.\deckid\par}
    \vspace{0.25cm}
    {\usebeamerfont{title}\color{RedTitle}\LARGE\bfseries \decktitletext\par}
    \vspace{0.7cm}
    {\large Zhigang Feng\par}
    \vspace{0.15cm}
    {\color{black!65}\deckdate\par}
    \vfill
    {\footnotesize\itshape\color{black!55}Build it on paper $\;\to\;$ solve it on the machine $\;\to\;$ push it to the frontier with AI\par}
  \end{frame}}

% ---------------------------------------------------------------------------
% Section divider (plain frame, not counted)
\newcommand{\sectiondivider}[2]{%
\begin{frame}[plain,noframenumbering]
\begin{center}
\vfill
{\Large\textbf{#1}\par}
\vspace{0.35cm}
{\normalsize #2\par}
\vfill
\end{center}
\end{frame}
}

% ---------------------------------------------------------------------------
% Boxes
\newtcolorbox{conceptbox}[1][]{colback=blue!3, colframe=RedTitle!55, boxrule=0.35mm,
  arc=1mm, left=2mm, right=2mm, top=1mm, bottom=1mm, #1}
\newtcolorbox{cautionbox}[1][]{colback=red!3, colframe=red!50!black, boxrule=0.35mm,
  arc=1mm, left=2mm, right=2mm, top=1mm, bottom=1mm, #1}
\newtcolorbox{examplebox}[1][]{colback=green!3, colframe=green!45!black, boxrule=0.35mm,
  arc=1mm, left=2mm, right=2mm, top=1mm, bottom=1mm, #1}
\newtcolorbox{takeawaybox}[1][]{colback=VeryLightGray, colframe=RedTitle!55, boxrule=0.35mm,
  arc=1mm, left=2mm, right=2mm, top=1mm, bottom=1mm, #1}
\newtcolorbox{researchbox}[1][]{enhanced, colback=violet!4, colframe=violet!60!black,
  boxrule=0.35mm, arc=1mm, left=2mm, right=2mm, top=1mm, bottom=1mm,
  fonttitle=\bfseries\small, title={Research in the wild}, #1}

\newcommand{\compacttables}{\renewcommand{\arraystretch}{1.15}}
\newcommand{\src}[1]{{\tiny\color{black!55}#1}}

% ---------------------------------------------------------------------------
%% Course map: the recurring "you are here" slide for ECON 282E.
%% Loaded by course_preamble.tex. Pattern follows AI-ECON-2026 lec01_map.tex, but the map
%% shows the whole ten-session course rather than one lecture.
%%
%%   \CourseMap{s}{label}   s = 1..10 highlights session s and prints "you are here: label"
%%                          (e.g. \CourseMap{1}{1.A2}); earlier sessions get a check mark.
%%   \CourseMap{0}{}        full overview, nothing highlighted.
%%
%% Note: TikZ option lists are comma-split before TeX conditionals run, so every \ifnum
%% below sits outside [...] option lists.

\newcommand{\CourseMap}[2]{%
\begin{center}
\begin{tikzpicture}[
    sess/.style={rectangle, rounded corners=2pt, draw=black!45, minimum width=1.34cm,
        minimum height=1.12cm, text width=1.26cm, align=center, inner sep=1pt},
    stage/.style={font=\tiny\bfseries, text=black!70},
]
  \foreach \i/\d/\t/\c in {%
      1/Sep 24/Intro \\ Growth/StageTrunk,
      2/Oct 1/HA $\cdot$ OLG \\ Policy/StageBranch,
      3/Oct 8/Num.\ DP \\ Python/StageMachine,
      4/Oct 15/Numerics \\ I/StageMachine,
      5/Oct 22/Numerics \\ II/StageMachine,
      6/Oct 29/The wall \\ \textbf{Midterm}/StageWall,
      7/Nov 5/ML \\ Deep nets/StageTools,
      8/Nov 12/RL \\ HA $+$ DL/StageTools,
      9/Nov 19/Text \\ LLMs/StageData,
      10/Dec 3/Agents \\ \textbf{Projects}/StageAgents} {
    \node[sess, fill=\c] (n\i) at ({1.46*(\i-1)}, 0)
      {{\scriptsize\bfseries S\i}\\[-1pt]{\tiny\color{black!60}\d}\\[1pt]{\tiny \t}};
    \ifnum\i<#1
      \node[font=\scriptsize, text=green!45!black] at ([xshift=-2.5mm,yshift=-2.2mm]n\i.north east) {$\checkmark$};
    \fi
    \ifnum\i=#1
      \draw[RedTitle, very thick, rounded corners=2pt]
        ([xshift=-0.6mm,yshift=-0.6mm]n\i.south west) rectangle ([xshift=0.6mm,yshift=0.6mm]n\i.north east);
      % keep the label inside the picture at both ends of the row
      \ifnum\i<3
        \node[font=\scriptsize\bfseries, text=RedTitle, anchor=south west] at ([yshift=1.2mm]n\i.north west)
          {$\blacktriangledown$\,you are here: #2};
      \else\ifnum\i>8
        \node[font=\scriptsize\bfseries, text=RedTitle, anchor=south east] at ([yshift=1.2mm]n\i.north east)
          {you are here: #2\,$\blacktriangledown$};
      \else
        \node[font=\scriptsize\bfseries, text=RedTitle, anchor=south] at ([yshift=1.2mm]n\i.north)
          {you are here: #2\,$\blacktriangledown$};
      \fi\fi
    \fi
  }
  % stage band under the sessions
  \foreach \a/\b/\lab/\c in {1/1/trunk/StageTrunk, 2/2/branches/StageBranch,
      3/5/the machine $\cdot$ classical methods/StageMachine, 6/6/the wall/StageWall,
      7/8/new tools/StageTools, 9/9/new data/StageData, 10/10/colleagues/StageAgents} {
    \fill[\c] ([yshift=-1.5mm]n\a.south west) rectangle ([yshift=-4.6mm]n\b.south east);
    \path ([yshift=-1.5mm]n\a.south west) -- ([yshift=-4.6mm]n\b.south east)
      node[midway, stage] {\lab};
  }
  \node[font=\footnotesize\itshape, text=black!70] at ([yshift=-9.5mm]$(n1.south)!0.5!(n10.south)$)
    {Build it on paper $\;\to\;$ solve it on the machine $\;\to\;$ push it to the frontier with AI};
\end{tikzpicture}
\end{center}
}


%%   \decktitle{1}{A1}{Who I Am \& Why This Course}{September 24, 2026}
%%   \begin{document}
%%   \makedecktitle
%%   ...
%%
%% Adapted from Courses/AI-ECON-2026/source/shared_preamble.tex (Metropolis theme,
%% concept/caution/example/takeaway boxes, \sectiondivider). Changes for this course:
%%   * course title block (\decktitle / \makedecktitle) instead of \makebeamertitle
%%   * \zh{...} typesets Chinese (book titles) under pdflatex via CJKutf8
%%   * researchbox for the "Research in the wild" frames
%%   * section pages off; TikZ libraries and the course map loaded here
%% Build with pdflatex only (two passes); tools/build_slides.py does this for you.

\documentclass[10pt,english,aspectratio=169]{beamer}

\usepackage[T1]{fontenc}
\usepackage{textcomp}
\usepackage[utf8]{inputenc}
\usepackage{babel}
\usepackage{CJKutf8}

\usepackage{amsmath, amssymb, amsthm, graphicx}
\usepackage{booktabs, multirow, tabularx}
\usepackage{tikz}
\usetikzlibrary{positioning,arrows.meta,shapes,calc,fit,backgrounds}
\usepackage{pgfplots}
\pgfplotsset{compat=1.16}
\usepackage[most]{tcolorbox}
\tcbuselibrary{skins, breakable}
\usepackage{listings}
\usepackage{xcolor}

\lstset{
  basicstyle=\ttfamily\small,
  keywordstyle=\color{blue!80!black}\bfseries,
  commentstyle=\color{green!50!black}\itshape,
  stringstyle=\color{orange!80!black},
  backgroundcolor=\color{gray!8},
  frame=single,
  rulecolor=\color{gray!40},
  breaklines=true,
  showstringspaces=false,
  numbers=none,
}

\usetheme{metropolis}
\usefonttheme{professionalfonts}
\metroset{sectionpage=none}

\definecolor{LightGray}{RGB}{230,230,230}
\definecolor{RedTitle}{RGB}{0,0,200}
\definecolor{VeryLightGray}{RGB}{245,245,245}
\definecolor{DarkText}{RGB}{0,0,0}
\definecolor{UCRBlue}{RGB}{0,45,106}
\definecolor{UCRGold}{RGB}{241,171,0}

% Stage colours of the course arc (used by the course map and elsewhere)
\colorlet{StageTrunk}{blue!12}
\colorlet{StageBranch}{cyan!14}
\colorlet{StageMachine}{gray!18}
\colorlet{StageWall}{red!14}
\colorlet{StageTools}{orange!20}
\colorlet{StageData}{green!16}
\colorlet{StageAgents}{violet!14}

\hypersetup{bookmarks=false,breaklinks=true,pdfborder={0 0 0},colorlinks=true,
  linkcolor=DarkText,urlcolor=RedTitle}

\setbeamercolor{background canvas}{bg=white}
\setbeamercolor{title}{fg=RedTitle}
\setbeamercolor{frametitle}{bg=VeryLightGray,fg=DarkText}
\setbeamercolor{normal text}{fg=DarkText}
\setbeamercolor{alerted text}{fg=RedTitle}
\setbeamersize{text margin left=5mm, text margin right=5mm}
\addtobeamertemplate{frametitle}{}{\vspace{-0.5em}}

\setbeamertemplate{navigation symbols}{}
\setbeamertemplate{footline}{%
  \hspace{5mm}{\usebeamerfont{page number in head/foot}\color{black!45}%
  ECON 282E \,$\cdot$\, \insertshortdeck}\hfill%
  \usebeamercolor[fg]{page number in head/foot}%
  \usebeamerfont{page number in head/foot}%
  \insertframenumber\,/\,\inserttotalframenumber\kern1em\vskip2pt%
}

% ---------------------------------------------------------------------------
% Chinese text under pdflatex (book titles etc.)
\newcommand{\zh}[1]{\begin{CJK*}{UTF8}{gbsn}#1\end{CJK*}}

% ---------------------------------------------------------------------------
% Deck title block
%   \decktitle{session}{deck id}{title}{date}
\newcommand{\insertshortdeck}{}
\newcommand{\decksession}{}
\newcommand{\deckid}{}
\newcommand{\decktitletext}{}
\newcommand{\deckdate}{}
\newcommand{\decktitle}[4]{%
  \renewcommand{\decksession}{#1}%
  \renewcommand{\deckid}{#2}%
  \renewcommand{\decktitletext}{#3}%
  \renewcommand{\deckdate}{#4}%
  \renewcommand{\insertshortdeck}{Session #1 \,$\cdot$\, Deck #1.#2}%
  \title{#3}%
  \author{Zhigang Feng}%
  \date{#4}%
}
\newcommand{\makedecktitle}{%
  \begin{frame}[plain,noframenumbering]
    \begin{tikzpicture}[remember picture,overlay]
      \fill[UCRBlue] (current page.north west) rectangle ([yshift=-0.55cm]current page.north east);
      \fill[UCRGold] ([yshift=-0.55cm]current page.north west) rectangle ([yshift=-0.62cm]current page.north east);
      \node[anchor=west, text=white, font=\small\bfseries] at ([xshift=5mm,yshift=-0.275cm]current page.north west)
        {ECON 282E \,$\cdot$\, Foundations of Macroeconomics};
      \node[anchor=east, text=white, font=\small] at ([xshift=-5mm,yshift=-0.275cm]current page.north east)
        {UC Riverside \,$\cdot$\, Fall 2026};
    \end{tikzpicture}
    \vspace*{1.1cm}
    {\usebeamerfont{title}\color{black!55}\large Session \decksession\ \,$\cdot$\, Deck \decksession.\deckid\par}
    \vspace{0.25cm}
    {\usebeamerfont{title}\color{RedTitle}\LARGE\bfseries \decktitletext\par}
    \vspace{0.7cm}
    {\large Zhigang Feng\par}
    \vspace{0.15cm}
    {\color{black!65}\deckdate\par}
    \vfill
    {\footnotesize\itshape\color{black!55}Build it on paper $\;\to\;$ solve it on the machine $\;\to\;$ push it to the frontier with AI\par}
  \end{frame}}

% ---------------------------------------------------------------------------
% Section divider (plain frame, not counted)
\newcommand{\sectiondivider}[2]{%
\begin{frame}[plain,noframenumbering]
\begin{center}
\vfill
{\Large\textbf{#1}\par}
\vspace{0.35cm}
{\normalsize #2\par}
\vfill
\end{center}
\end{frame}
}

% ---------------------------------------------------------------------------
% Boxes
\newtcolorbox{conceptbox}[1][]{colback=blue!3, colframe=RedTitle!55, boxrule=0.35mm,
  arc=1mm, left=2mm, right=2mm, top=1mm, bottom=1mm, #1}
\newtcolorbox{cautionbox}[1][]{colback=red!3, colframe=red!50!black, boxrule=0.35mm,
  arc=1mm, left=2mm, right=2mm, top=1mm, bottom=1mm, #1}
\newtcolorbox{examplebox}[1][]{colback=green!3, colframe=green!45!black, boxrule=0.35mm,
  arc=1mm, left=2mm, right=2mm, top=1mm, bottom=1mm, #1}
\newtcolorbox{takeawaybox}[1][]{colback=VeryLightGray, colframe=RedTitle!55, boxrule=0.35mm,
  arc=1mm, left=2mm, right=2mm, top=1mm, bottom=1mm, #1}
\newtcolorbox{researchbox}[1][]{enhanced, colback=violet!4, colframe=violet!60!black,
  boxrule=0.35mm, arc=1mm, left=2mm, right=2mm, top=1mm, bottom=1mm,
  fonttitle=\bfseries\small, title={Research in the wild}, #1}

\newcommand{\compacttables}{\renewcommand{\arraystretch}{1.15}}
\newcommand{\src}[1]{{\tiny\color{black!55}#1}}

% ---------------------------------------------------------------------------
%% Course map: the recurring "you are here" slide for ECON 282E.
%% Loaded by course_preamble.tex. Pattern follows AI-ECON-2026 lec01_map.tex, but the map
%% shows the whole ten-session course rather than one lecture.
%%
%%   \CourseMap{s}{label}   s = 1..10 highlights session s and prints "you are here: label"
%%                          (e.g. \CourseMap{1}{1.A2}); earlier sessions get a check mark.
%%   \CourseMap{0}{}        full overview, nothing highlighted.
%%
%% Note: TikZ option lists are comma-split before TeX conditionals run, so every \ifnum
%% below sits outside [...] option lists.

\newcommand{\CourseMap}[2]{%
\begin{center}
\begin{tikzpicture}[
    sess/.style={rectangle, rounded corners=2pt, draw=black!45, minimum width=1.34cm,
        minimum height=1.12cm, text width=1.26cm, align=center, inner sep=1pt},
    stage/.style={font=\tiny\bfseries, text=black!70},
]
  \foreach \i/\d/\t/\c in {%
      1/Sep 24/Intro \\ Growth/StageTrunk,
      2/Oct 1/HA $\cdot$ OLG \\ Policy/StageBranch,
      3/Oct 8/Num.\ DP \\ Python/StageMachine,
      4/Oct 15/Numerics \\ I/StageMachine,
      5/Oct 22/Numerics \\ II/StageMachine,
      6/Oct 29/The wall \\ \textbf{Midterm}/StageWall,
      7/Nov 5/ML \\ Deep nets/StageTools,
      8/Nov 12/RL \\ HA $+$ DL/StageTools,
      9/Nov 19/Text \\ LLMs/StageData,
      10/Dec 3/Agents \\ \textbf{Projects}/StageAgents} {
    \node[sess, fill=\c] (n\i) at ({1.46*(\i-1)}, 0)
      {{\scriptsize\bfseries S\i}\\[-1pt]{\tiny\color{black!60}\d}\\[1pt]{\tiny \t}};
    \ifnum\i<#1
      \node[font=\scriptsize, text=green!45!black] at ([xshift=-2.5mm,yshift=-2.2mm]n\i.north east) {$\checkmark$};
    \fi
    \ifnum\i=#1
      \draw[RedTitle, very thick, rounded corners=2pt]
        ([xshift=-0.6mm,yshift=-0.6mm]n\i.south west) rectangle ([xshift=0.6mm,yshift=0.6mm]n\i.north east);
      % keep the label inside the picture at both ends of the row
      \ifnum\i<3
        \node[font=\scriptsize\bfseries, text=RedTitle, anchor=south west] at ([yshift=1.2mm]n\i.north west)
          {$\blacktriangledown$\,you are here: #2};
      \else\ifnum\i>8
        \node[font=\scriptsize\bfseries, text=RedTitle, anchor=south east] at ([yshift=1.2mm]n\i.north east)
          {you are here: #2\,$\blacktriangledown$};
      \else
        \node[font=\scriptsize\bfseries, text=RedTitle, anchor=south] at ([yshift=1.2mm]n\i.north)
          {you are here: #2\,$\blacktriangledown$};
      \fi\fi
    \fi
  }
  % stage band under the sessions
  \foreach \a/\b/\lab/\c in {1/1/trunk/StageTrunk, 2/2/branches/StageBranch,
      3/5/the machine $\cdot$ classical methods/StageMachine, 6/6/the wall/StageWall,
      7/8/new tools/StageTools, 9/9/new data/StageData, 10/10/colleagues/StageAgents} {
    \fill[\c] ([yshift=-1.5mm]n\a.south west) rectangle ([yshift=-4.6mm]n\b.south east);
    \path ([yshift=-1.5mm]n\a.south west) -- ([yshift=-4.6mm]n\b.south east)
      node[midway, stage] {\lab};
  }
  \node[font=\footnotesize\itshape, text=black!70] at ([yshift=-9.5mm]$(n1.south)!0.5!(n10.south)$)
    {Build it on paper $\;\to\;$ solve it on the machine $\;\to\;$ push it to the frontier with AI};
\end{tikzpicture}
\end{center}
}


\graphicspath{{./figures/}}

\lstset{language=Python, basicstyle=\ttfamily\scriptsize, frame=none,
        backgroundcolor=\color{gray!8}, aboveskip=2pt, belowskip=2pt}

\decktitle{3}{B1}{Your Toolkit: Python, Environments, and Version Control}{Thursday, October 8, 2026}

\begin{document}

\makedecktitle

%=============================================================================
\begin{frame}{Where We Are}
\CourseMap{3}{3.B1}
\vspace{-1mm}
\begin{center}
\small \textbf{Block B:} \textbf{3.B1} your toolkit $\;\cdot\;$ \textbf{3.B2} from loops to arrays $\;\cdot\;$ \textbf{3.B3} PyTorch, and coding with AI $\;\cdot\;$ \textbf{3.B4} the data workflow.
\end{center}
\end{frame}

%=============================================================================
\begin{frame}{Four Error Messages You Are Going to Meet}
\renewcommand{\arraystretch}{1.25}
\begin{center}\footnotesize
\begin{tabularx}{\linewidth}{@{}>{\raggedright\arraybackslash}p{5.0cm} >{\raggedright\arraybackslash}X >{\raggedright\arraybackslash}p{2.6cm}@{}}
\toprule
\textbf{What you will see} & \textbf{What it means} & \textbf{Fixed by} \\
\midrule
\texttt{ModuleNotFoundError: no module named 'numpy'} & the package is not installed \emph{in the interpreter you are using} & environments \\
\texttt{command not found: python} & the shell cannot find an interpreter at all & installation, \texttt{PATH} \\
\texttt{AttributeError} on a function that exists in the docs & you are running a different version than the docs describe & environments \\
``It worked yesterday.'' & you changed something and cannot say what & version control \\
\bottomrule
\end{tabularx}
\end{center}
\vspace{1mm}
\begin{conceptbox}\small
Not one of these is a Python problem. They are all \textbf{environment} problems, and they will cost you more hours this quarter than any algorithm in Block A. This deck is the twenty minutes that prevents them.
\end{conceptbox}
\end{frame}

%=============================================================================
\begin{frame}{How the Pieces Fit Together}
\begin{columns}[c]
\begin{column}{0.54\textwidth}
\centering
\begin{tikzpicture}[font=\scriptsize,
  lyr/.style={draw=black!45, rounded corners=2pt, minimum width=6.6cm, minimum height=0.78cm, align=center}]
  \node[lyr, fill=blue!6,   draw=RedTitle] (a) at (0,2.55) {\textbf{Editor / notebook}\\[-1pt]{\scriptsize JupyterLab, VS Code, Cursor}};
  \node[lyr, fill=green!6]  (b) at (0,1.50) {\textbf{Packages}\\[-1pt]{\scriptsize NumPy, SciPy, pandas, matplotlib, PyTorch}};
  \node[lyr, fill=orange!6] (c) at (0,0.45) {\textbf{Python interpreter}\\[-1pt]{\scriptsize inside a named environment}};
  \node[lyr, fill=black!5]  (d) at (0,-0.60) {\textbf{Operating system}\\[-1pt]{\scriptsize macOS, Windows, Linux --- or a compute node}};
  \draw[->, black!50] (d) -- (c); \draw[->, black!50] (c) -- (b); \draw[->, black!50] (b) -- (a);
\end{tikzpicture}
\end{column}
\begin{column}{0.43\textwidth}
\begin{conceptbox}[title=\textbf{The distinction that matters}]\small
The editor is \emph{not} the thing that runs your code. It sends text to an interpreter, and which interpreter it picks is a setting --- usually the one you forgot to check.
\end{conceptbox}
\medskip
\begin{examplebox}\footnotesize
Every ``it works for you but not for me'' in this course is a disagreement about one of the middle two layers.
\end{examplebox}
\end{column}
\end{columns}
\vspace{1mm}
\src{Adapted from QM Lec.~3, ``How the Pieces Fit Together''.}
\end{frame}

%=============================================================================
\begin{frame}{Where You Will Actually Type}
\begin{columns}[T]
\begin{column}{0.49\textwidth}
\small
\begin{center}\footnotesize
\begin{tabularx}{\linewidth}{@{}>{\raggedright\arraybackslash}p{1.8cm} >{\raggedright\arraybackslash}X@{}}
\toprule
\textbf{Tool} & \textbf{Best for} \\
\midrule
JupyterLab & exploring, plotting, homework you hand in \\
VS Code / Cursor & multi-file projects, debugging, AI assistance inline \\
PyCharm & large codebases, heavy refactoring \\
\bottomrule
\end{tabularx}
\end{center}
\medskip
\footnotesize
\textbf{For this course:} notebooks for the labs and homework, an editor once your project outgrows one file --- which will happen in Session 5.
\end{column}
\begin{column}{0.48\textwidth}
\begin{cautionbox}[title=\textbf{A notebook is not a program}]\footnotesize
Cells can be run out of order, so a notebook that looks correct may depend on a variable you have since deleted.\\[1mm]
Before you submit anything: \textbf{Restart Kernel and Run All}. If it fails, it was never working.
\end{cautionbox}
\medskip
\begin{examplebox}\footnotesize
The same discipline is why the reference numbers in Block A were produced by a script submitted to a compute node, not by a notebook someone had been clicking around in.
\end{examplebox}
\end{column}
\end{columns}
\vspace{1mm}
\src{Adapted from QM Lec.~3, ``Jupyter, VS Code/Cursor, and PyCharm''.}
\end{frame}

%=============================================================================
\begin{frame}{Language, Library, Model --- Three Different Things}
\begin{columns}[T]
\begin{column}{0.52\textwidth}
\small
\begin{itemize}\footnotesize
  \item \textbf{Python} is the language: syntax, control flow, functions, classes. It is slow, and that is fine, because it is glue.
  \item \textbf{NumPy, SciPy, PyTorch} are libraries: compiled code with a Python interface. The arithmetic happens in C or CUDA, not in your loop.
  \item \textbf{An AI assistant} is a model behind an interface. It writes Python; it does not run it, and it does not know whether the answer is right.
\end{itemize}
\end{column}
\begin{column}{0.45\textwidth}
\begin{takeawaybox}[title=\textbf{Why this matters for speed}]\small
Python being slow is only a problem when \emph{Python} does the arithmetic. Push the loop into the library and the language stops mattering --- which is the whole of 3.B2.
\end{takeawaybox}
\medskip
\begin{examplebox}\footnotesize
And why it matters for trust: the library is tested by thousands of users; the assistant's output has been tested by nobody.
\end{examplebox}
\end{column}
\end{columns}
\vspace{1mm}
\src{Adapted from QM Lec.~3, ``Python, PyTorch, and AI: How They Relate''.}
\end{frame}

%=============================================================================
\begin{frame}[fragile]{Installing: Who Does What}
\begin{columns}[T]
\begin{column}{0.50\textwidth}
\small
\textbf{Anaconda / Miniconda} bundles three separate things people confuse: a Python interpreter, a package manager, and an environment manager. Miniconda is the same thing without the 3\,GB of packages you will not use.
\medskip

\textbf{An editor is not a distribution.} Installing VS Code does not install Python; it goes looking for one.
\medskip
\footnotesize
On Windows, tick \emph{``Add to PATH''}. Almost every \texttt{command not found} traces back to that box.
\end{column}
\begin{column}{0.47\textwidth}
\begin{lstlisting}[language=bash]
$ python --version
Python 3.11.9

$ which python
/opt/miniconda3/envs/econ282e/bin/python
\end{lstlisting}
\vspace{-1mm}
\begin{conceptbox}\footnotesize
Run both of these \emph{before} asking anyone for help. They answer ``which Python am I actually using?'', which is the first question every answer depends on.
\end{conceptbox}
\end{column}
\end{columns}
\vspace{1mm}
\src{Adapted from QM Lec.~3, ``Anaconda, Python, and Editors'' and ``Installing Python''.}
\end{frame}

%=============================================================================
\begin{frame}[fragile]{Environments: One Habit, Most of the Breakage Gone}
\begin{columns}[T]
\begin{column}{0.50\textwidth}
\small
An environment is a private folder holding one interpreter and one set of package versions. Projects stop fighting over versions because they never share any.
\medskip
\begin{lstlisting}[language=bash]
conda create -n econ282e python=3.11
conda activate econ282e
pip install numpy scipy pandas \
    matplotlib jupyter torch tqdm
pip freeze > requirements.txt
\end{lstlisting}
\end{column}
\begin{column}{0.47\textwidth}
\begin{cautionbox}[title=\textbf{Two rules}]\footnotesize
\textbf{1.} Packages are installed \emph{into} an environment. Installing while the wrong one is active is the single most common cause of \texttt{ModuleNotFoundError}.\\[1mm]
\textbf{2.} Prefer \texttt{conda} \emph{or} \texttt{pip} within one environment; mixing them for the same package leads to states neither tool can repair.
\end{cautionbox}
\medskip
\begin{takeawaybox}\footnotesize
\texttt{requirements.txt} is what makes your results reproducible by someone else --- including the version of you who returns in March.
\end{takeawaybox}
\end{column}
\end{columns}
\vspace{1mm}
\src{The \texttt{install} commands follow QM Lec.~3, ``Managing Packages''; environment creation is new --- the source covers installing packages but never creating an environment to put them in.}
\end{frame}

%=============================================================================
\begin{frame}[fragile]{Git: So That ``It Worked Yesterday'' Is a Question You Can Answer}
\begin{columns}[T]
\begin{column}{0.48\textwidth}
\small
Git records \emph{snapshots} of a folder with a message attached. That is enough to answer the three questions that otherwise cost an afternoon: what changed, when, and what did it look like before.
\medskip
\begin{lstlisting}[language=bash]
git init
git add vfi.py
git commit -m "VFI matches Brock-Mirman"
git log --oneline
git diff HEAD~1
\end{lstlisting}
\end{column}
\begin{column}{0.49\textwidth}
\begin{conceptbox}[title=\textbf{Commit when a number changes}]\small
The useful habit is not daily commits --- it is committing at the moment a result becomes reproducible, with the result in the message.
\end{conceptbox}
\medskip
\begin{cautionbox}\footnotesize
\textbf{Never commit} data you do not own, model checkpoints, or anything with a password in it. A \texttt{.gitignore} is written once per project and saves you from all three.
\end{cautionbox}
\medskip
\footnotesize
Homework in this course is submitted as a PDF, the code that produced it, and an \texttt{AI\_LOG.md}. Git is how you make sure those three agree.
\end{column}
\end{columns}
\vspace{1mm}
\src{Written for this course; version control does not appear anywhere in QM Lec.~3.}
\end{frame}

%=============================================================================
\begin{frame}{When You Need a Bigger Machine}
\begin{columns}[T]
\begin{column}{0.51\textwidth}
\small
Nothing in Block A needed more than a laptop: the entire morning's computation ran in under a second. That stops being true in Sessions 7 and 8, when training a network replaces iterating an operator.
\medskip
\begin{itemize}\footnotesize
  \item \textbf{Google Colab} gives a free GPU in a browser. Good for a lab, bad for anything you need to reproduce --- the machine is gone when you close the tab.
  \item \textbf{A cluster} gives you a scheduler. You write a job script, submit it, and collect the output; the login node is for editing, never for computing.
\end{itemize}
\end{column}
\begin{column}{0.46\textwidth}
\begin{takeawaybox}[title=\textbf{The portable habit}]\small
Put the computation in a \textbf{script} that takes its settings from the top of the file and writes its results to disk. Then the same file runs on a laptop, in Colab, or under a scheduler, unchanged.
\end{takeawaybox}
\medskip
\begin{examplebox}\footnotesize
That is exactly how this morning's numbers were produced: one script, one output file, and every figure drawn from it.
\end{examplebox}
\end{column}
\end{columns}
\vspace{1mm}
\src{Written for this course; Colab is mentioned in QM Lec.~2's programming-environment frame but not covered.}
\end{frame}

%=============================================================================
\begin{frame}[fragile]{The Four Atomic Types}
\begin{columns}[T]
\begin{column}{0.50\textwidth}
\begin{lstlisting}
year      = 2026          # int
inflation = 0.031         # float
country   = "Chile"       # str
recession = False         # bool

type(inflation)           # <class 'float'>
int("2026") + 1           # 2027
float("3.1") * 2          # 6.2
\end{lstlisting}
\end{column}
\begin{column}{0.47\textwidth}
\small
Python infers the type; you do not declare it. That is convenient and occasionally dangerous.
\medskip
\begin{cautionbox}\footnotesize
\texttt{1/2} is \texttt{0.5} in Python 3 but \texttt{0} in Python 2 and in most of C. Integer division is written \texttt{//} --- and is exactly how you index into a grid.
\end{cautionbox}
\medskip
\begin{examplebox}\footnotesize
A float here is the float64 of 3.A1: about sixteen significant digits, and gaps that widen as the numbers grow.
\end{examplebox}
\end{column}
\end{columns}
\vspace{1mm}
\src{Adapted from QM Lec.~3, ``Atomic Data Types''.}
\end{frame}

%=============================================================================
\begin{frame}[fragile]{Lists, Slicing, and Text}
\begin{columns}[T]
\begin{column}{0.52\textwidth}
\begin{lstlisting}
g = [2.1, 2.9, -0.4, 1.8, 2.5]

g[0]        # 2.1   first
g[-1]       # 2.5   last
g[1:3]      # [2.9, -0.4]
g[:2], g[3:]
len(g), sum(g) / len(g)

g.append(3.0)       # modifies g
sorted(g)[:2]       # two worst years

name = "  gdp Growth "
name.strip().title()   # 'Gdp Growth'
f"inflation was {0.031:.1%}"
\end{lstlisting}
\end{column}
\begin{column}{0.45\textwidth}
\small
\textbf{Slices are half-open}: \texttt{g[1:3]} gives elements 1 and 2. It reads badly and it is the right convention --- \texttt{g[:i] + g[i:]} reconstructs the list for every \texttt{i}.
\medskip
\begin{conceptbox}\footnotesize
Indexing from \textbf{zero} and slicing half-open are the two conventions behind every off-by-one error you will write this quarter. They are also what make \texttt{k[i]}, \texttt{k[i+1]} in an interpolation formula line up.
\end{conceptbox}
\end{column}
\end{columns}
\vspace{1mm}
\src{Adapted from QM Lec.~3, ``Lists: Ordered Collections'' and ``Working with Text''.}
\end{frame}

%=============================================================================
\begin{frame}[fragile]{The One That Bites: Reference or Copy?}
\begin{columns}[T]
\begin{column}{0.50\textwidth}
\begin{lstlisting}
a = [1, 2, 3]
b = a            # NOT a copy
b[0] = 99
a                # [99, 2, 3]

c = a.copy()     # a real copy
c[0] = 1
a                # [99, 2, 3]
\end{lstlisting}
\vspace{-1mm}
\small
\texttt{b = a} makes a second \emph{name} for one object. Assignment never copies.
\end{column}
\begin{column}{0.47\textwidth}
\begin{cautionbox}[title=\textbf{Why this will bite you}]\footnotesize
The same rule governs NumPy arrays and PyTorch tensors --- and slices of them are \textbf{views}, not copies.\\[1mm]
In a value-function loop, \texttt{V\_new = V} followed by an update silently destroys the very thing you wanted to compare against.
\end{cautionbox}
\medskip
\begin{examplebox}\footnotesize
That is why 3.A2's loop writes \texttt{V, n = Vn, n + 1} --- binding a \emph{new} array --- and never writes into the old one.
\end{examplebox}
\end{column}
\end{columns}
\vspace{1mm}
\src{Adapted from QM Lec.~3, ``Warning: Reference vs Copy''.}
\end{frame}

%=============================================================================
\begin{frame}[fragile]{Dictionaries, Tuples, Sets}
\begin{columns}[T]
\begin{column}{0.51\textwidth}
\begin{lstlisting}
par = {"alpha": 0.36, "beta": 0.95,
       "delta": 1.0}
par["beta"]              # 0.95
par["sigma"] = 0.007     # add a key
for k, v in par.items():
    print(k, v)

bounds = (0.05, 0.60)    # tuple: fixed
# bounds[0] = 0.1        -> TypeError

set([1, 1, 2, 3])        # {1, 2, 3}
\end{lstlisting}
\end{column}
\begin{column}{0.46\textwidth}
\small
\begin{itemize}\footnotesize
  \item \textbf{Dictionaries} are how you carry a calibration around. One object, named fields, and it serialises straight to JSON --- which is how a results file records the parameters that produced it.
  \item \textbf{Tuples} are immutable. Use them for things that must not change: grid bounds, an array shape, a pair of coordinates.
  \item \textbf{Sets} deduplicate and test membership fast.
\end{itemize}
\end{column}
\end{columns}
\vspace{1mm}
\src{Adapted from QM Lec.~3, ``Dictionaries: Key-Value Data'' and ``Tuples and Sets''.}
\end{frame}

%=============================================================================
\begin{frame}{Takeaways}
\begin{takeawaybox}
\begin{itemize}\small
  \item The errors that cost you time this quarter will be \textbf{environment} errors, not Python errors. One environment per project, created on the first day, removes most of them.
  \item \texttt{which python} and \texttt{python -\,-version} answer the question every other answer depends on. Run them before asking for help.
  \item The editor does not run your code and a notebook is not a program. \textbf{Restart and Run All} before you submit anything.
  \item Git is how ``it worked yesterday'' becomes a question with an answer. Commit when a number changes, and put the number in the message.
  \item Assignment never copies. \texttt{b = a} gives one object two names --- and the same rule governs NumPy arrays and PyTorch tensors, where slices are views.
\end{itemize}
\end{takeawaybox}
\end{frame}

%=============================================================================
\begin{frame}{Bridge: The Algorithm Is on the Board. Put It in an Array.}
\begin{columns}[T]
\begin{column}{0.53\textwidth}
\small
You now have a machine that runs Python and a project that can be reproduced. What you do not have is the one idea that separates code which finishes from code which does not.
\medskip

This morning's Bellman update is a double loop: for every state, for every choice. Written that way in Python it is unusably slow. Written as a single array expression it is the same arithmetic, the same answer to the last bit --- and it runs about \textbf{65 times faster}.
\end{column}
\begin{column}{0.44\textwidth}
\begin{conceptbox}[title=\textbf{Next (3.B2)}]\small
Control flow, functions and classes --- then NumPy, broadcasting, and the vectorised Bellman update.\\[1mm]
We rebuild this morning's solver as a \texttt{GrowthModel} class, and measure what each change bought.
\end{conceptbox}
\end{column}
\end{columns}
\end{frame}

\end{document}
